\section*{Sources, comparisons, and verification record}
\addcontentsline{toc}{section}{Sources, comparisons, and verification record}
We have treated the attached original \cite{AC26} as the source of its
own formulas. We extracted its text, compared the relevant formulas
with the supplied transcription, and inspected its original page~66
for the orientation computation used in D.1.4. The algebra, maps,
finite quotient models and compatibility calculations used above
are given in our proofs and in the accompanying executed programs.
The workbench, course and survey helped locate claims; none of their
status labels enters a proof.

Our use of \cite{E26} includes the statement and proof of Lemma~2.1,
the reference linearization in Corollary~2.2, the geometric quotient
in Section~2, the finite modifications in Section~3, and Sections~7--8.
The passage from its linearization to our periods is B.1.1 and B.1.4.
The height input is also computed in B.1.3 and applied explicitly
in D.1.2. B.2.3 locates a global-parameter descent correction to
Remark~8.1. E.1.5 identifies the remaining simultaneous degeneration
needed for Conjecture~8.2; no statement about the topology of an
unconstructed compact central fibre is made.

For the deformation comparison we read the fixed-parameter
classification argument in \cite[Section~6.2]{Y26} and the precise
non-isotriviality statement in \cite[Proposition~4.7]{VL26}.
A.4.3 gives the maps between these methods and the present
calculation. The additional starting surfaces and changed period
families in \cite{VTwo26,VCombs26,LR26,W26} were checked through
their introductions and stated main theorems. The lattice in A.2
belongs to the displayed $(3,4,\infty)$ period system. Substituting
a different starting surface changes those matrices, and the
proofs above do not perform that substitution. The coprime
$(3k,4l)$ substitution that we do perform is B.3.1. This states
the exact scope of the classification through the actual period
matrices and finite norm maps.

The zero of $\wp$ in C.1 is derived from the explicit section and
the trigonometric cusp expansion. We read
\cite[Theorem~1]{DI08} and its discussion of Eichler--Zagier,
but did not use an unopened original paper as a premise.
For the description of line bundles and their sections in D, we
use the Appell--Humbert correspondence and the Riemann criterion
in \cite[Theorems~4.1 and 4.4]{Schnell19}; the substitution of
the actual period lattice and the rational-kernel quotient are
written in D.1 and D.2. E.2.2 states the exact gluing and
characteristic-number reduction to which the theorems of
\cite{Cerf71,Smale61,EK62} are applied.

We searched the public literature on 12~November~2026 using the
paper titles and identifiers above and combinations of
``complex six-sphere'', ``classification'', ``period parameter'',
``Klingler'', ``monodromy'', ``Engel Conjecture 8.2'' and
``degeneration''. The report records the query groups and inspected
sources. This search also located an important version change:
\cite[Theorems~1.3--1.4]{DGKX26v1} stated miniversality, the
$2\Z$ identification and the full automorphism group; the current
version \cite{DGKX26v2} omits those deformation and classification
sections. We record the two texts separately and infer no reason
for that revision. We then checked the actual root-curve, cusp
cohomology and logarithmic-tangent arguments in the historical
version. A.5.2 gives their explicit derivation for the present
charts, including the coarse invariant exponents and the normal
quotient at every cusp stratum. The resulting tangent dimensions
and our Kodaira--Spencer calculation give the stated miniversal
disc. The complete audit and independently executed weighted
cohomology checks are recorded in the report. We claim no priority
for any result.

All decimal tables identify working precision and independently
computed expressions. Exact lattice, Smith-form and polynomial
identities use integer or rational arithmetic. The scripts,
their actual outputs, the item-by-item report and a statement
ledger accompany this paper. Exploratory floating-point grids
are identified as such; a certified enclosure is used only
where its error bounds and completed computation are supplied.
